% 第 5 章
\bichapter{布线与布线后优化}{Routing and Postroute Optimization}
\label{chap:route}

本章介绍 IC Compiler II 路由器 Zroute 的布线能力，以及布线后优化功能。Zroute 面向多核硬件架构，可高效处理 45~nm 及以下先进工艺设计规则与可制造性设计（DFM）任务。

本章包括以下主题：
\begin{itemize}
  \item Zroute 简介（Introduction to Zroute）
  \item 基本 Zroute 流程（Basic Zroute Flow）
  \item 布线前提条件（Prerequisites for Routing）
  \item 定义 Via（Defining Vias）
  \item 插入 Via Ladder（Inserting Via Ladders）
  \item 检查可布线性（Checking Routability）
  \item 布线约束（Routing Constraints）
  \item 布线应用选项（Routing Application Options）
  \item 时钟网络布线（Routing Clock Nets）
  \item 关键网络布线（Routing Critical Nets）
  \item 次级电源/地 pin 布线（Routing Secondary Power and Ground Pins）
  \item 信号网络布线（Routing Signal Nets）
  \item 网络屏蔽（Shielding Nets）
  \item 执行布线后优化（Performing Postroute Optimization）
  \item 分析并修复信号电迁移违例（Analyzing and Fixing Signal Electromigration Violations）
  \item 执行 ECO 布线（Performing ECO Routing）
  \item 在 GUI 中布线（Routing Nets in the GUI）
  \item 清理已布线网络（Cleaning Up Routed Nets）
  \item 分析布线结果（Analyzing the Routing Results）
  \item 保存布线信息（Saving Route Information）
  \item 推导掩模颜色（Deriving Mask Colors）
  \item 插入与移除 Cut Metal 形状（Inserting and Removing Cut Metal Shapes）
\end{itemize}

% ============================================================
\section{Zroute 简介}
\subsection*{Introduction to Zroute}

Zroute 包含五个布线引擎：全局布线（global routing）、轨道分配（track assignment）、详细布线（detail routing）、ECO 布线与布线验证。全局布线、轨道分配与详细布线可通过任务专用命令或自动布线命令调用；ECO 布线与布线验证则使用任务专用命令。

Zroute 主要特性包括：
\begin{itemize}
  \item 全布线步骤（含全局布线、轨道分配、详细布线）支持多核多线程；
  \item 真实连通模型：矩形相触即视为电气连通，无需导线中心线重合；
  \item 动态迷宫网格：可 off-grid 连接 pin，同时保留网格路由器的速度优势；
  \item 多边形管理器：识别多边形，并理解 DRC 面向多边形；
  \item 详细布线阶段并发优化设计规则、天线规则、导线与 via；
  \item 详细布线阶段并发插入冗余 via；
  \item 全局/轨道/详细布线内建 soft rule 支持；
  \item 时序与串扰驱动的全局布线、轨道分配、详细布线与 ECO 布线；
  \item 智能设计规则处理（合并冗余违例、智能收敛）；
  \item 带层约束与非默认布线规则（NDR）的 net group 布线；
  \item 时钟布线；
  \item 布线验证；
  \item 以 soft rule 方式优化 DFM 与可制造性（DFY）；
  \item 支持先进设计规则，如多重图案化（multiple patterning）。
\end{itemize}

% ============================================================
\section{基本 Zroute 流程}
\subsection*{Basic Zroute Flow}

图~\ref{fig:basic-zroute-flow} 给出基本 Zroute 流程，包括时钟布线、信号布线、DFM 优化与布线验证。

\begin{figure}[htbp]
\centering
\fbox{\parbox{0.85\textwidth}{\centering\vspace{1.5cm}
\textit{（原书 Figure 63：Basic Zroute Flow）}\\[0.3em]
\small 请参阅原版 PDF 插图；可将截图放入 \texttt{figures/} 目录后用 \texttt{\textbackslash includegraphics} 替换。
\vspace{1.5cm}}}
\caption{基本 Zroute 流程}
\label{fig:basic-zroute-flow}
\end{figure}

% ============================================================
\section{布线前提条件}
\subsection*{Prerequisites for Routing}

运行 Zroute 前，block 与物理库须满足：
\begin{itemize}
  \item \textbf{库要求}：Zroute 从 technology file 获取全部设计规则信息，布线前须确保规则已在 technology file 中定义。详见 \textit{Synopsys Technology File and Routing Rules Reference Manual}。
  \item \textbf{Block 要求}：
  \begin{itemize}
    \item 电源/地网络已在设计规划后、布局前布线完成（见 \textit{IC Compiler II Design Planning User Guide}）；
    \item 已完成时钟树综合与优化（见第~\ref{chap:cts}~章）；
    \item 估计拥塞可接受；
    \item 估计时序可接受（约 0~ns slack）；
    \item 估计最大电容与 transition 无违例。
  \end{itemize}
\end{itemize}
后三项可通过 \cmd{check\_routability} 验证布局可布线性（见「检查可布线性」）。

% ============================================================
\section{定义 Via}
\subsection*{Defining Vias}

路由器支持以下 via 定义类型：
\begin{itemize}
  \item \textbf{Simple via 与 simple via array}：simple via 为单 cut via，由 cut 层及 cut/金属层矩形的高宽指定；simple via array 为多 cut 的 simple via。可用于：带 NDR 的时钟/信号布线、冗余 via 插入、带先进 via 规则的 PG 布线。
  \item \textbf{Custom via}：由任意 Manhattan 多边形集合构成的多 cut、异形 via；\textbf{仅}可用于冗余 via 插入。
\end{itemize}

Via 定义来源：
\begin{itemize}
  \item 设计库关联 technology file 的 \texttt{ContactCode} 段；
  \item LEF 文件中的 via rule \texttt{GENERATE} 语句（定义 simple via array）；
  \item 用户自定义 via（\cmd{create\_via\_def}）。
\end{itemize}

若希望 Zroute 将某 via 定义用于信号布线，须满足：\appopt{is\_default} 为 \cmd{true}，且 \appopt{is\_excluded\_for\_signal\_routing} 为 \cmd{false}。此外，Zroute 还会使用 technology file fat via table 或 \cmd{create\_routing\_rule} 所定义 NDR 中显式指定的非默认 via。

\subsection{从 LEF 读取 Via 定义}
\subsubsection*{Reading Via Definitions from a LEF File}

使用 \cmd{read\_tech\_lef} 读取 LEF 中 via rule \texttt{GENERATE} 语句。支持如下 LEF 语法：
\begin{lstlisting}
via_rule_name GENERATE [DEFAULT]
   LAYER lower_layer_name
      ENCLOSURE lower_overhang1 lower_overhang2
   LAYER upper_layer_name
      ENCLOSURE upper_overhang1 upper_overhang2
   LAYER cut_layer_name
      RECT llx lly urx ury
      SPACING x_spacing BY y_spacing
\end{lstlisting}
不支持 \texttt{WIDTH} 与 \texttt{RESISTANCE} 语句。

\subsection{创建 Via 定义}
\subsubsection*{Creating a Via Definition}

使用 \cmd{create\_via\_def} 创建 via 定义；创建后可在设计中任意处使用（类似 technology file 的 \texttt{ContactCode}）。用户 via 定义保存在设计库中，仅该设计可用。可定义 simple via、simple via array 与 custom via。

\subsubsection{定义 Simple Via}
\subsubsection*{Defining Simple Vias}

\begin{lstlisting}
create_via_def
   -cut_layer layer
   -cut_size {width height}
   -upper_enclosure {width height}
   -lower_enclosure {width height}
   [-min_rows number_of_rows]
   [-min_columns number_of_columns]
   [-min_cut_spacing distance]
   [-cut_pattern cut_pattern]
   [-is_default]
   [-force]
   via_def_name
\end{lstlisting}

无需指定 enclosure 层：工具根据 technology file 中与 via 层相邻的金属层自动确定。默认创建单 cut via；定义 via array 时指定 \opt{-min\_rows}/\opt{-min\_columns} 与 \opt{-min\_cut\_spacing}。默认 cut pattern 为满阵列；可用 \opt{-cut\_pattern} 修改。覆盖已有定义须用 \opt{-force}。

\begin{lstlisting}
icc2_shell> create_via_def design_via1_HV \
  -cut_layer VIA12 -cut_size {0.05 0.05} \
  -lower_enclosure {0.02 0.0} -upper_enclosure {0.0 0.2}

icc2_shell> create_via_def design_via12 -cut_pattern "01 10"
\end{lstlisting}

用 \cmd{report\_via\_defs} 报告用户 via 定义。

\subsubsection{定义 Custom Via}
\subsubsection*{Defining Custom Vias}

\begin{lstlisting}
create_via_def
   -shapes { {layer {coordinates} [mask_constraint]} ... }
   [-lower_mask_pattern alternating | uniform]
   [-upper_mask_pattern alternating | uniform]
   [-force]
   via_def_name
\end{lstlisting}

\opt{-shapes} 须为一个 via 层与两个金属 enclosure 层指定形状；每层可有多个形状。

\begin{lstlisting}
icc2_shell> create_via_def design_H_shape_via \
   -shapes { {VIA12 {0.035 -0.035} {0.100 0.035}}
             {VIA12 {-0.100 -0.035} {-0.035 0.035}}
             {METAL1 {-0.130 -0.035} {0.130 0.035}}
             {METAL2 {-0.100 -0.065} {-0.035 0.065}}
             {METAL2 {0.035 -0.065} {0.100 0.065}}
             {METAL2 {-0.100 -0.035} {0.100 0.035}} }
\end{lstlisting}

双图案化时，可对形状或 enclosure 层分配 mask constraint，二者不可同时使用：
\begin{itemize}
  \item 对形状：在 \opt{-shapes} 中用 \cmd{mask\_one}/\cmd{mask\_two}/\cmd{mask\_three}；
  \item 对 enclosure 层：用 \opt{-lower\_mask\_pattern}/\opt{-upper\_mask\_pattern}（\cmd{uniform} 或 \cmd{alternating}），创建 via 时再指定首个形状的 mask。
\end{itemize}

\begin{lstlisting}
icc2_shell> create_via_def VIA12 \
   -shapes { {M1 {-0.037 -0.010} {-0.017 0.01} mask_two}
             {M1 {0.017 -0.010} {0.037 0.01} mask_one}
             {VIA1 {-0.037 -0.01} {0.037 0.01} mask_two}
             {M2 {-0.044 -0.010} {0.044 0.010} mask_one} }
icc2_shell> create_via -net n1 -via_def VIA12 \
   -origin {100.100 200.320}
\end{lstlisting}
